% Target ~350 words.

\paragraph{Acoustic feature engineering.} Early automatic dysarthria
severity work relied on hand-crafted acoustic features: long-term-average
spectra, perturbation measures (jitter, shimmer), and rhythm
metrics~\cite{Tu2017XVectorDysarthria,Berisha2014AcousticMarkers}. MFCC
+ SVR/Ridge remains the default baseline. These features are
interpretable but brittle to channel and recording variability and require
per-speaker calibration in deployment.

\paragraph{Self-supervised representations.} A second wave of work
fine-tunes large SSL backbones (wav2vec~2.0~\cite{Baevski2020wav2vec2},
HuBERT~\cite{Hsu2021HuBERT}, WavLM~\cite{Chen2022WavLM}) on dysarthric
classification or regression
heads~\cite{Hernandez2020SSLDysarthria,Joshi2021IndianDysarthria,Schu2023SSLDys}.
Fine-tuning yields strong in-distribution accuracy but destroys the
pre-trained geometry, requires labels at training time, and tends to
overfit speaker-specific spectral idiosyncrasies (we observe this directly
in our LOSO comparison: MFCC trajectory DTW achieves the highest global
$\rho$ but the worst LOSO $\rho$).

\paragraph{Frozen SSL probes.} Recent perceptual-quality work demonstrates
that frozen SSL representations can be probed with simple distance metrics
and recover useful speech-quality
signals~\cite{Millet2022SSLPerceptual,Janbakhshi2021PCADys}. Our
contribution is the trajectory-geometric framing: rather than mean-pooling
or training a classifier head, we treat the per-frame embedding sequence
as a curve in a manifold and measure deviation from a reference curve via
DTW. To our knowledge, this is the first dysarthria severity metric that
operates at frame-level trajectory geometry in a frozen SSL latent space
without any task-specific training.

\paragraph{Trajectory-based methods in atypical speech.} Dynamic time
warping has long been used in dysarthric ASR to handle the 2--3$\times$
slower production
rate~\cite{Bhat2022DTWDysarthria,Salvador2007FastDTW}. We borrow the
alignment machinery and length-normalisation but apply it to learned SSL
representations rather than acoustic features, and we use it for severity
estimation rather than recognition.

\paragraph{Cross-disorder and cross-language extension.} Severity
estimation has been pursued separately for Parkinson's
disease~\cite{Vasquez-Correa2019CNNPD,Pompili2020SSLPD} and ALS, with
limited cross-disorder portability. Our cross-corpus extension (Section~5)
tests whether the SSL articulatory manifold generalises across disorders
(PD, dysarthria), languages (English, Spanish, Italian), and recording
conditions, using a shared reference basis with per-corpus normalisation.
